\clearpage\fancyhead[L]{\small GAUSS AN BODE.}\hypertarget{native-356-lower}{}\label{native-356-lower}
\begin{center}\textbf{GAUSS AN BODE.}\par\vspace{4pt}\hrule\vspace{4pt}\fontsize{10}{12}\selectfont Astronomisches Jahrbuch für 1815. Seite 190\qquad194. Berlin 1812.\par\vspace{4pt}\hrule\vspace{4pt}\end{center}

\begin{flushright}Göttingen den 5. Mai 1812.\end{flushright}
\par\indent Der Wunsch, von Ihnen, eben so wie im vorigen Jahre, auch in diesem einige Pallasbeobachtungen um die Zeit der Opposition (10. Juni) zu erhalten, veranlasst mich, Sie hierauf besonders aufmerksam zu machen, und was ich in dieser Beziehung bisher beobachtet habe, mitzutheilen. Am 4. April habe ich diesen Planeten zum erstenmal wieder gesehen, aber erst am 9ten beobachtet. Die gestrige war eben so wie die Beobachtung vom 9. April für Declination nicht vortheilhaft, daher ich diese ganz weglasse. Der Planet hat nicht viel mehr als die 10te Grösse.

\begin{center}
\fontsize{10}{12}\selectfont\setlength{\tabcolsep}{3pt}\renewcommand{\arraystretch}{1.15}\begin{tabular}{c|c|c|c}
1812 & Mittlere Zeit & AR. $\PallasSymbol$ & Abw. Nordl. \\
\hline
April\ \ 9 & $11^{\mathrm u}\;32^{\mathrm m}\;27^{\mathrm s}$ & $169^\circ\;5'\;58''4$ & \\
Mai\ \ 2 & $10\;36\;6$ & $268\;42\;5.3$ & $21^\circ\;4'\;37''7$ \\
4 & $10\;24\;16$ & $268\;32\;27.4$ & \\
\end{tabular}
\end{center}
\SourceLeaf{357}\par\indent Hr. \textsc{Nicolai} findet aus der Vergleichung mit den Elementen, nach welchen er seine Ephemeride berechnet hat, folgende Unterschiede:

\begin{center}
\fontsize{10}{12}\selectfont\setlength{\tabcolsep}{3pt}\renewcommand{\arraystretch}{1.15}\begin{tabular}{c|c|c}
 & AR. & Decl. \\
\hline
April\ \ 9 & $+30''6$ & \\
Mai\ \ 2 & $+57.6$ & $-5''4$ \\
4 & $+49.4$ & \\
\end{tabular}
\end{center}

\par\indent Meine Beobachtungen halte ich für gut, die Vergleichung geschah mit 93, 96, 95 Herkules nach \textsc{Piazzi}'s Catalog. Den Unterschied in der ger. Aufst. bin ich geneigt den Störungen vom Saturn und Mars zuzuschreiben, die bei meinen Rechnungen noch vernachlässigt sind, die ich aber bald auch vornehmen werde. Die Anzahl aller periodischen von mir berechneten Störungen, die von Jupiter herrühren, beläuft sich auf 400, wobei alle weggelassen sind, die unter $1''$ betragen. Mehrere merkwürdige Resultate haben sich bei dieser Gelegenheit dargeboten, die ich aber vorerst noch zurückhalten muss. Denn meine ganze Rechnung kann nur als eine vorläufige angesehen werden, die noch einmal und in noch grösserer Ausdehnung wiederholt werden muss.

\par\indent Von den andern neuen Planeten habe ich noch keine beobachtet. Blos die Gegend, wo die Juno steht, habe ich gestern Abend in sehr geringer Höhe, wo die kleinen Sterne noch schlecht zu sehen waren, ins Fernrohr gebracht, wo es mir schien, als wenn der Platz, wo Juno nach der Ephemeride stehen sollte, leer, hingegen unter einer etwa $25'$ kleinern AR. ein Stern 9ter bis 10ter Grösse befindlich war, welches vielleicht die Juno gewesen ist. Dies aber nur zur vorläufigen Notiz, was sich bald näher bestimmen lassen wird.

\par\indent Ich theile Ihnen bei dieser Gelegenheit noch einige andere astronomische Beobachtungen mit. — Von der Erscheinung des 2ten Cometen wurde ich am 9ten Dec. v. J. durch Hrn. v. \textsc{Lindenau} benachrichtigt, und an demselben Tage fand ich ihn sogleich. Hier meine sämmtlichen Beobachtungen desselben, die ich alle für gut halte, die letzte Declination ausgenommen, die nicht ganz so zuverlässig ist.

\begin{center}
\fontsize{10}{12}\selectfont\setlength{\tabcolsep}{3pt}\renewcommand{\arraystretch}{1.15}\begin{tabular}{c|c|c|c}
 & Mittlere Zeit & AR. & Decl. \\
\hline
1811 Dec.\ \ 9 & $10^{\mathrm u}\;6^{\mathrm m}\;52^{\mathrm s}$ & $63^\circ\;49'\;41''4$ & $10^\circ\;21'\;55''5$ S. \\
11 & $10\;34\;1$ & $63\;33\;18.0$ & $8\;39\;46.4$ \\
12 & $\;8\;5\;52$ & $63\;26\;25.3$ & $7\;54\;25.9$ \\
1812 Jan.\ \ 3 & $\;6\;54\;46$ & $62\;37\;6.4$ & $9\;47\;14.7$ N. \\
4 & $\;6\;54\;15$ & $62\;40\;38.6$ & $10\;29\;57.8$ \\
Febr.\ \ 2 & $11\;51\;52$ & $68\;9\;17.9$ & $26\;1\;34$ \\
\end{tabular}
\end{center}

\Needspace{15\baselineskip}\par\indent Gleich nach meiner Beobachtung vom 12. Dec. berechnete ich die parabolischen Elemente, die ich indessen weglasse, um statt ihrer die vom Hrn. \textsc{Nicolai} verbesserten zu geben, welche dieser sehr geschickte Rechner gleich nach meiner Beobachtung vom 4. Januar herausgebracht hat.

\begin{center}\fontsize{10}{12}\selectfont\setlength{\tabcolsep}{2pt}\renewcommand{\arraystretch}{1.15}\begin{tabular}{lll}
\multicolumn{3}{l}{Perihel Zeit Nov 11.39211 Meridian von Göttingen}\\
— Länge & $.\;.\;.\;.\;.\;.\;.\;.\;.\;.\;.\;.\;.\;.\;.$ & $47^\circ\;39''\;36''1$\\
— Logarithm. des Abstandes & $.\;.\;.\;.\;.\;.\;.\;.$ & $0.2011781$\\
Länge des $\AscendingNode$ & $.\;.\;.\;.\;.\;.\;.\;.\;.\;.\;.\;.\;.\;.\;.\;.$ & $92^\circ\;53'\;44''2$\\
Neigung der Bahn & $.\;.\;.\;.\;.\;.\;.\;.\;.\;.\;.\;.\;.\;.$ & $31\;32\;38.7$\\
\multicolumn{3}{l}{Bewegung rechtläufig.}\\
\end{tabular}\end{center}

\Needspace{10\baselineskip}\par\indent Die Vergleichung dieser Elemente mit meinen Beobachtungen gab nach Hrn. \textsc{Nicolai}'s Rechnung folgende Resultate.
\SourceLeaf{358}\begin{center}
Differenz
\end{center}

\begin{center}
\fontsize{10}{12}\selectfont\setlength{\tabcolsep}{3pt}\renewcommand{\arraystretch}{1.15}\begin{tabular}{c|c|c||c|c|c}
1811 & AR. & Decl. & 1812 & AR. & Decl. \\
\hline
Dec. 9 & $-5''0$ & $+3''4$ & Jan. 3 & $-6''5$ & $+0''3$ \\
11 & $+2.9$ & $-32.7$ & 4 & $+2.8$ & $-12.7$ \\
12 & $+8.5$ & $-27.6$ & Febr. 2 & $-113.7$ & $+18.1$ \\
\end{tabular}
\end{center}

\par\indent Hr. \textsc{Nicolai} will nächstens seine Elemente noch einmal nach den letzten im Februar gemachten
Beobachtungen verbessern.

\Needspace{30\baselineskip}\par\indent Hier noch einige Beobachtungen von Sternbedeckungen:

\begin{center}\fontsize{9}{11}\selectfont\setlength{\tabcolsep}{2pt}\renewcommand{\arraystretch}{1.15}\begin{tabular}{llrrrll}
1811 Octob. 23 & Stern 6r Grösse & $7^{\mathrm u}$ & $21^{\mathrm m}$ & $40^{\mathrm s}1$ & Eintritt & Harding\\
— & 7r & 8 & 14 & 40.4 & — & Harding\\
— & 6.\ 7r & 8 & 30 & 33.8 & — & Gauss\\
24 & 7.\ 8r & 7 & 38 & 48.8 & — & Harding\\
1812 Febr. 1 & $\eta$ Virginis & 12 & 21 & 26.7 & — & Gauss\\
19 & $\gamma$ Tauri & 7 & 22 & 15.0 & Austritt & Gerling\\
— & 71 Tauri & 9 & 56 & 3.4 & Eintritt & Gerling\\
 &  &  &  & 3.8 & — & Harding\\
 &  &  &  & 3.9 & — & Gauss\\
— & $1\,\theta$ Tauri & 10 & 58 & 56.9 & — & Harding\\
 &  &  &  & 57.4 & — & Gauss\\
 &  &  &  & 57.9 & — & Gerling\\
 &  & 11 & 55 & 34.6 & Austritt & Gauss\\
 &  &  &  & 38.6 & — & Harding\\
— & $2\,\theta$ Tauri & 10 & 58 & 12.0 & Eintritt & Gauss und Harding\\
 &  & 11 & 59 & 6.6 & Austritt & Gauss\\
 &  &  &  & 8.6 & Austritt & Harding\\
— & 160 Mayer & 12 & 3 & 39.5 & Eintritt & Harding\\
 &  &  &  & 40.0 & — & Gauss\\
— & 162 Mayer & 12 & 26 & 27.0 & — & Gauss\\
 &  &  &  & 27.5 & — & Harding\\
20 & 111 Tauri & 11 & 7 & 46.8 & Eintritt & Gerling\\
— & Anonyma & 11 & 36 & 54.0 & — & Harding\\
— & 117 Tauri & 13 & 0 & 25.0 & — & Gerling\\
 &  &  &  & 25.1 & — & Harding\\
\end{tabular}\end{center}

\par\indent Die Bedeckung von $\eta\;\text{\AstroSymbols ♍}$ am 23. April fand hier in Göttingen gar nicht statt.

\par\indent Noch habe ich eine Beobachtung von der Bedeckung $\lambda\;\text{\AstroSymbols ♒}$ vom 2. Sept. 1811 nachzuholen, die
mir Hr. Obrist v. \textsc{Tischleder} mit dem Auftrage, sie Ihnen mitzutheilen, zugesandt hat.

\par\indent Feste Rosenberg ob Cronach: Eintr. $\lambda$ Aquarii 1811 Sept. 2 $10^{\mathrm u}$ $25^{\mathrm m}$ $5^{\mathrm s}4$ M. Z.

\par\indent Hr. \textsc{Wachter} hat hieraus und aus den Beobachtungen von Göttingen und Mannheim folgende
Conjunctionszeiten berechnet: Mannheim $10^{\mathrm u}$ $16^{\mathrm m}$ $38^{\mathrm s}01$, Göttingen $10^{\mathrm u}$ $22^{\mathrm m}$ $29^{\mathrm s}69$, Cronach $10^{\mathrm u}$ $28^{\mathrm m}$ $3^{\mathrm s}50$.

\par\indent Es ist also die Länge von Cronach aus der Vergleichung mit

\begin{center}
\fontsize{10}{12}\selectfont\setlength{\tabcolsep}{3pt}\renewcommand{\arraystretch}{1.15}\begin{tabular}{ll}
Göttingen & $35^{\mathrm m}$ $55^{\mathrm s}81$ \\
Mannheim & 35 \quad $57.49$ \\
im Mittel & 35 \quad $56.65$ oder in Graden $28^\circ$ $59'$ $10''$ \\
\end{tabular}
\end{center}
\SourceLeaf{359}\par\indent Die Polhöhe setzt der Obrist $=50^\circ$ $14'$ $42''6$.

\par\indent An unserer neuen Sternwarte wird auch diesen Sommer thätig gearbeitet. Nächstens erwarte ich
einen 12zölligen Multiplicationskreis v. \textsc{Reichenbach}. Hr. Dr. \textsc{Gerling} kommt als Lehrer der Mathematik
an das Lyceum in Cassel. Er wird seine Stelle im September antreten. Seine Geschicklichkeit und
sein Eifer machen es wünschenswerth, dass er auch dort noch sich astronomischen Beschäftigungen möge
unterziehen können. Seine Probeschrift über Sternbedeckungen und Finsternisse und besonders über die
bei uns (aber nicht in Berlin) ringförmige Sonnenfinsterniss von 1820 (7. Sept.) wird jetzt gedruckt.
